\chapter{Zaključek}
\label{ch:zakljucek}

V diplomskem delu smo zasnovali in razvili spletno in mobilno aplikacijo za upravljanje delovnega časa in kadrovskih procesov, ki v enotnem sistemu združuje evidenco delovnega časa, razporejanje izmen z možnostjo samodejnega generiranja osnutka urnika in upravljanje odsotnosti, ob tem pa podpira dostop na osnovi vlog za administratorja, kadrovsko službo, vodjo in zaposlenega ter nadzorni vmesnik za upravljanje več organizacij na eni platformi. Dosegli smo vse cilje, opredeljene v razdelku~\ref{sec:cilji}: zajeli smo zahteve po vlogah, zasnovali podatkovni model in arhitekturo (poglavje~\ref{ch:nacrtovanje}), razvito rešitev opisali po modulih (poglavje~\ref{ch:sistem}) in jo ovrednotili s testnimi uporabniki (poglavje~\ref{ch:evalvacija}).

Glede na rezultate evalvacije hipotezo H3, da bodo uporabniki aplikacijo ocenili kot uporabno in enostavno za uporabo, potrdimo: povprečna ocena SUS na vzorcu 12 testnih uporabnikov je znašala 88,1, kar po pridevniški lestvici Bangorja in sodelavcev~\cite{bangor2009determining} sodi v razred ``odlično''. Hipotezo H2 o časovnem prihranku pri samodejnem razporejanju lahko na podlagi izvedene evalvacije potrdimo le delno, saj časovnega prihranka v primerjavi z ročnim razporejanjem nismo izmerili neposredno, temveč jo posredno podpira hitra in nemotena izvedba naloge razporejanja med testiranjem. Hipotezo H1, da enoten sistem zmanjša potrebo po ročnem usklajevanju podatkov med ločenimi orodji, neposredno nismo merili s testnimi uporabniki, temveč jo utemeljujemo z zasnovo samega sistema (poglavje~\ref{ch:nacrtovanje}), v katerem so evidenca delovnega časa, razporejanje izmen in odsotnosti od začetka povezani na skupnem podatkovnem modelu, ne da bi bilo med njimi potrebno ročno prenašanje podatkov. K podobnemu sklepu vodijo tudi raziskave s področja digitalizacije kadrovskih procesov: Mahmoud in sodelavci~\cite{mahmoud2025digitalhrm} na vzorcu zaposlenih pokažejo, da digitalni kadrovski informacijski sistemi statistično značilno zmanjšajo zapletenost nalog in potrebo po ročnem delu ter s tem izboljšajo učinkovitost kadrovske funkcije.

\section{Omejitve rešitve}

Razvita rešitev ima kljub doseženim ciljem več omejitev. Reševalnik za samodejno razporejanje je požrešna hevristika, ki izmene dodeljuje kronološko in ne poišče nujno globalno najboljše razporeditve za celotno obdobje, kot smo pojasnili v razdelkih~\ref{sec:razporejanje-osebja} in~\ref{sec:sistem-razporejanje}. Pri zahtevnejših urnikih z veliko omejitvami bi eksaktni ali metahevristični pristop lahko dal boljši rezultat, a na račun daljšega časa izvajanja in zahtevnejšega vzdrževanja kode. Evalvacija je potekala na majhnem, priložnostnem vzorcu udeležencev (kolegi, družina, prijatelji), zato dobljenih ocen SUS ne moremo posplošiti na širšo populacijo uporabnikov kadrovskih sistemov, vendar pa je potrebno povedati, da so udeleženci sodili v ciljno skupino uporabnikov. Sistem prav tako še ne podpira obračuna plač niti integracije z zunanjimi kadrovskimi ali računovodskimi sistemi, ki jih ponujajo nekatere komercialne rešitve iz razdelka~\ref{sec:komercialne-resitve}.

\section{Nadaljnje delo}

Rešitev bi lahko nadgradili v več smereh. Reševalnik razporejanja bi lahko razširili z možnostjo, da zaposleni za posamezno izmeno oddajo eksplicitno ponudbo namesto zgolj splošne preference, kar bi vodji dalo močnejši signal pri odločanju, ali pa reševalnik nadomestili z metahevrističnim oziroma eksaktnim pristopom, na primer s celoštevilskim programiranjem, ki bi lahko poleg zasedenosti optimiziral tudi strošek dela. Smiselna razširitev bi bila tudi podpora obračunu plač na podlagi evidentiranega delovnega časa. Evalvacijo bi bilo koristno ponoviti na večjem in bolj raznolikem vzorcu uporabnikov, morda v obliki daljšega uporabniškega preizkusa v realnem podjetju, kar bi omogočilo tudi merjenje dejanskega prihranka časa pri razporejanju izmen, ne le subjektivne ocene uporabnikov.
